\documentclass[10pt,a4paper]{article}
\usepackage[T1]{fontenc}
\usepackage[utf8]{inputenc}
\usepackage{lmodern}
\usepackage[a4paper,top=1.6cm,bottom=1.6cm,left=2.1cm,right=2.1cm]{geometry}
\usepackage{microtype}
\usepackage{enumitem}
\usepackage[hidelinks]{hyperref}
\setlength{\parindent}{0pt}
\setlength{\parskip}{0.55em}
\pagestyle{empty}
\newcommand{\letterhead}{%
\begin{flushright}
\textbf{Dr.\ Kundai Farai Sachikonye}\\
AIMe Registry for Artificial Intelligence\\
Email: \href{mailto:kundai.sachikonye@bitspark.com}{kundai.sachikonye@bitspark.com}\\
Phone: +49 162 6386344
\end{flushright}
\vspace{0.4em}
5 October 2026\par
\vspace{0.4em}
The Editor\\
Editorial Office\par}
\newcommand{\declarations}{%
\textbf{Declarations.}
\begin{itemize}[leftmargin=1.2em,itemsep=0.15em,topsep=0.2em]
\item The manuscript is original, has not been published previously, and is not under consideration for publication elsewhere.
\item I am the sole author, have approved the submitted version, and take full responsibility for its content.
\item I declare no competing financial or non-financial interests.
\item This work received no specific grant from any funding agency in the public, commercial, or not-for-profit sectors.
\end{itemize}}
\newcommand{\signoff}{%
Thank you for considering this manuscript. I look forward to the reviewers' assessment and am happy to supply any further information the editor requires.

Yours sincerely,\par
\vspace{0.8em}
Dr.\ Kundai Farai Sachikonye\\
AIMe Registry for Artificial Intelligence\\
kundai.sachikonye@bitspark.com \quad|\quad +49 162 6386344}

\begin{document}
\letterhead

\textbf{Re: Submission of manuscript, ``On Neuro-Musculoskeletal System Continuous Charge Distribution Dynamics: Continuous Charge Subtraction in Human Orthogonal Behavioural Conditions''}

Dear Editor,

I am pleased to submit the above manuscript for consideration as an original research article in your journal.

Whole-body calorimetry reports energy expenditure in watts, but it cannot separate the charge redistribution that supports voluntary thought from the charge redistribution that supports voluntary movement: both draw on overlapping neural substrates, and both are present in every realistic recording. The manuscript addresses this measurement problem directly. It shows that four ordinary behavioural conditions (deep sleep, REM sleep, meditative single-intent locomotion, and seated wakefulness) form an orthogonal basis for the charge-redistribution subspace of the body, treated as a closed, non-grounded bioelectrical circuit. Subtracting between these conditions isolates component charge rates for movement, thought, perception, and housekeeping, expressed in coulombs per second rather than watts.

The principal contributions are:
\begin{itemize}[leftmargin=1.2em,itemsep=0.25em]
\item \textbf{A subtraction protocol} that converts metabolic power to charge rate via $Q=\sqrt{2CU}$ on per-second charge--discharge cycles, combining the Attwell--Laughlin split of brain metabolism with a cardiac reference for the perceptual component.
\item \textbf{A stability criterion}, the activity--sleep mirror law $0.8 \le \mathcal{C}/\mathcal{E} \le 1.2$, which identifies the days on which continuous subtraction is numerically stable.
\item \textbf{An independent internal check}: the REM dream charge is predicted rather than fitted, and the measured dream-to-thought charge ratio of 0.975 lies within 2.5\% of the predicted unity.
\item \textbf{Validation} on 86 nights of hypnogram data and four instrumented runs, yielding $Q_{\mathrm{motor}} = 290.9$~mC/s, $Q_{\mathrm{thought}} = 133.5$~mC/s, $Q_{\mathrm{perception}} = 70.9$~mC/s and $Q_{\mathrm{dream}} = 130.2$~mC/s, with run-to-run variation of $\pm 3.1\%$ across five seed splits of the longitudinal record.
\end{itemize}

The procedure requires only consumer-grade wearables (photoplethysmography, accelerometry, and sleep-stage scoring) and a single in-clinic electrocardiogram, so it can be deployed immediately in longitudinal human phenotyping. The manuscript closes by deriving clinical signatures for Parkinson's disease, the Alzheimer's prodrome, and anaesthesia depth from the component-wise charge trajectories. I believe the work will interest readers in physiology, neuroscience, wearable sensing, and clinical phenotyping.

\declarations

\signoff
\end{document}
